\documentclass[10pt,a4paper]{amsart}
\usepackage{amsmath,amssymb,amsthm,graphicx}
% Retained minimal preamble for the Stacks Project chapter more-algebra.tex (More on Algebra).
% The chapter file pinned-original-source.tex carries only \input{preamble} on its first line, so no
% preamble text lives inside the chapter itself.  The Stacks Project uses one single shared preamble for
% every chapter; the snapshot stages it identically as sources/stacks-chow/preamble.tex and
% sources/stacks-etale-cohomology/preamble.tex.  The block below is copied mechanically, line for line,
% from sources/stacks-chow/preamble.tex lines 14-18 (the Xy-pic setup) and 154-217 (the native theorem
% declarations together with the mathematical macros).  Nothing was added, removed or redefined; the
% wrapper class and the amsmath/amssymb/amsthm/graphicx base packages are supplied by the assembler.
% Chapter: more-algebra.tex -- More on Algebra.
% For commutative diagrams we use Xy-pic
\usepackage[all]{xy}

% We use 2cell for 2-commutative diagrams.
\xyoption{2cell}

% Theorem environments.
%
\theoremstyle{plain}
\newtheorem{theorem}[subsection]{Theorem}
\newtheorem{proposition}[subsection]{Proposition}
\newtheorem{lemma}[subsection]{Lemma}

\theoremstyle{definition}
\newtheorem{definition}[subsection]{Definition}
\newtheorem{example}[subsection]{Example}
\newtheorem{exercise}[subsection]{Exercise}
\newtheorem{situation}[subsection]{Situation}

\theoremstyle{remark}
\newtheorem{remark}[subsection]{Remark}
\newtheorem{remarks}[subsection]{Remarks}

\numberwithin{equation}{subsection}

% Macros
%
\def\lim{\mathop{\mathrm{lim}}\nolimits}
\def\colim{\mathop{\mathrm{colim}}\nolimits}
\def\Spec{\mathop{\mathrm{Spec}}}
\def\Hom{\mathop{\mathrm{Hom}}\nolimits}
\def\Ext{\mathop{\mathrm{Ext}}\nolimits}
\def\SheafHom{\mathop{\mathcal{H}\!\mathit{om}}\nolimits}
\def\SheafExt{\mathop{\mathcal{E}\!\mathit{xt}}\nolimits}
\def\Sch{\mathit{Sch}}
\def\Mor{\mathop{\mathrm{Mor}}\nolimits}
\def\Ob{\mathop{\mathrm{Ob}}\nolimits}
\def\Sh{\mathop{\mathit{Sh}}\nolimits}
\def\NL{\mathop{N\!L}\nolimits}
\def\CH{\mathop{\mathrm{CH}}\nolimits}
\def\proetale{{pro\text{-}\acute{e}tale}}
\def\etale{{\acute{e}tale}}
\def\QCoh{\mathit{QCoh}}
\def\Ker{\mathop{\mathrm{Ker}}}
\def\Im{\mathop{\mathrm{Im}}}
\def\Coker{\mathop{\mathrm{Coker}}}
\def\Coim{\mathop{\mathrm{Coim}}}

% Boxtimes
%
\DeclareMathSymbol{\boxtimes}{\mathbin}{AMSa}{"02}

%
% Macros for moduli stacks/spaces
%
\def\QCohstack{\mathcal{QC}\!\mathit{oh}}
\def\Cohstack{\mathcal{C}\!\mathit{oh}}
\def\Spacesstack{\mathcal{S}\!\mathit{paces}}
\def\Quotfunctor{\mathrm{Quot}}
\def\Hilbfunctor{\mathrm{Hilb}}
\def\Curvesstack{\mathcal{C}\!\mathit{urves}}
\def\Polarizedstack{\mathcal{P}\!\mathit{olarized}}
\def\Complexesstack{\mathcal{C}\!\mathit{omplexes}}
% \Pic is the operator that assigns to X its picard group, usage \Pic(X)
% \Picardstack_{X/B} denotes the Picard stack of X over B
% \Picardfunctor_{X/B} denotes the Picard functor of X over B
\def\Pic{\mathop{\mathrm{Pic}}\nolimits}
\def\Picardstack{\mathcal{P}\!\mathit{ic}}
\def\Picardfunctor{\mathrm{Pic}}
\def\Deformationcategory{\mathcal{D}\!\mathit{ef}}

% --- appended 2026-10-12 (Chief): the project own definitions of the reference/slogan/history
% display environments, copied mechanically from the shared Stacks preamble lines 8-12
% (\usepackage{verbatim} plus the three \newenvironment lines, which the chapter preamble keeps
% verbatim from the project preamble). Nothing added, removed or redefined: these environments
% carry no mathematical content and the project itself typesets them as comments. Without these
% three lines a counted statement whose own text uses them fails with Environment slogan undefined.
\usepackage{verbatim}
\newenvironment{reference}{\comment}{\endcomment}
\newenvironment{slogan}{\comment}{\endcomment}
\newenvironment{history}{\comment}{\endcomment}
\usepackage{cleveref}
\crefname{originallemma}{original lemma}{original lemmas}
\Crefname{originallemma}{Original lemma}{Original lemmas}
\makeatletter
\newlabel{more-algebra-lemma-Koszul-regular-flat-locally-regular}{{068Q}{source}{}{sourcepointer-1}{}}
\newlabel{more-algebra-lemma-Koszul-regular-flat-locally-regular@cref}{{[originallemma][0][]068Q}{[1][0][]source}{}{}{}}
\makeatother
\crefname{originallemma}{original lemma}{original lemmas}
\Crefname{originallemma}{Original lemma}{Original lemmas}
\makeatletter
\newlabel{lemma-make-nonzero-divisor}{{068P}{source}{}{sourcepointer-2}{}}
\newlabel{lemma-make-nonzero-divisor@cref}{{[originallemma][0][]068P}{[1][0][]source}{}{}{}}
\makeatother
\crefname{originallemma}{original lemma}{original lemmas}
\Crefname{originallemma}{Original lemma}{Original lemmas}
\makeatletter
\newlabel{lemma-koszul-regular-flat-base-change}{{062H}{source}{}{sourcepointer-3}{}}
\newlabel{lemma-koszul-regular-flat-base-change@cref}{{[originallemma][0][]062H}{[1][0][]source}{}{}{}}
\makeatother
\crefname{originallemma}{original lemma}{original lemmas}
\Crefname{originallemma}{Original lemma}{Original lemmas}
\makeatletter
\newlabel{lemma-independence-of-generators}{{066A}{source}{}{sourcepointer-4}{}}
\newlabel{lemma-independence-of-generators@cref}{{[originallemma][0][]066A}{[1][0][]source}{}{}{}}
\makeatother
\crefname{originallemma}{original lemma}{original lemmas}
\Crefname{originallemma}{Original lemma}{Original lemmas}
\makeatletter
\newlabel{lemma-truncate-koszul-regular}{{068M}{source}{}{sourcepointer-5}{}}
\newlabel{lemma-truncate-koszul-regular@cref}{{[originallemma][0][]068M}{[1][0][]source}{}{}{}}
\makeatother
\title{Making a Koszul regular sequence regular after a faithfully flat smooth base change}
\author{The Stacks Project Authors}
\date{stacks-project:more-algebra.tex:local-snapshot-20261012; GNU FDL 1.2 or later}
\begin{document}
\maketitle
\setcounter{subsection}{1}
\noindent\textit{Editorial note.} All delivered mathematical text is the author's own, copied from the pinned Stacks Project chapter file more-algebra.tex; no statement or proof step was rewritten, reordered or removed.\par\medskip
\noindent\textit{Editorial note.} The chapter file pinned-original-source.tex carries only \textbackslash{}input\{preamble\} on its first line, so no preamble text lives inside the chapter. The retained file sources/stacks-more-algebra/retained-preamble.tex is a mechanical, line-for-line copy of the project's single shared preamble as staged in the snapshot (Xy-pic setup plus theorem declarations and mathematical macros). Nothing was added, removed or redefined, and the excerpt serves only as preamble, never as a span coordinate.\par\medskip
\noindent\textit{Editorial note.} Both counted spans start at the \textbackslash{}begin line of their environment, so the statement and the author's proof are delivered whole; surrounding prose of the section is deliberately outside the counted spans and is carried as standard prerequisites instead.\par\medskip
\noindent\textit{Editorial note.} The citation fields (source\_id, source\_version, authors, source\_url, license.name, source\_title) are identical to those of every other card this lane has produced from this chapter, so the batch selector treats them as one source.\par\medskip
\noindent\textit{Editorial note.} Benign-screen note: the proof contains the phrase 'it is clear that' / 'is obvious' once, used as a legitimate step inside a complete argument; the author proves the statement, so this is not a hint-only proof.\par\medskip
\noindent\textit{Allowed standard prerequisites (editorial).}\begin{itemize}
\item The kernel of multiplication by the sequence is zero for a Koszul regular sequence. used at proof lines 7827-7830. Provides the vanishing input for the nonzerodivisor argument.
\item Criterion making one element a nonzerodivisor after adjoining variables. Lemma \textbackslash{}ref\{lemma-make-nonzero-divisor\}, cited at proof lines 7831-7833. Makes the new first element a nonzerodivisor.
\item Koszul regularity under flat base change and under changing generators, and truncation of Koszul regular sequences. Lemma \textbackslash{}ref\{lemma-koszul-regular-flat-base-change\}, Lemma \textbackslash{}ref\{lemma-independence-of-generators\}, and Lemma \textbackslash{}ref\{lemma-truncate-koszul-regular\}, cited at proof lines 7834-7839. Give the induction step.
\end{itemize}\medskip
\ifdefined\hypertarget\hypertarget{sourcepointer-1}{}\fi
\noindent\textit{Source reference.} Original lemma 068Q: focused-seeds/AGP/sources/stacks-snapshot-20261011/tags.txt, source lines 8014--8014. The lemma's own canonical tag, quoted so that readers can resolve the label inside the delivered text; it points at the counted statement itself and carries no proof content.\par\medskip
\ifdefined\hypertarget\hypertarget{sourcepointer-2}{}\fi
\noindent\textit{Source reference.} Original lemma 068P: focused-seeds/AGP/sources/stacks-snapshot-20261011/tags.txt, source lines 8013--8013. Navigation pointer to a statement the counted proof refers to but does not contain; nothing load-bearing is delegated to it.\par\medskip
\ifdefined\hypertarget\hypertarget{sourcepointer-3}{}\fi
\noindent\textit{Source reference.} Original lemma 062H: focused-seeds/AGP/sources/stacks-snapshot-20261011/tags.txt, source lines 7796--7796. Navigation pointer to a statement the counted proof refers to but does not contain; nothing load-bearing is delegated to it.\par\medskip
\ifdefined\hypertarget\hypertarget{sourcepointer-4}{}\fi
\noindent\textit{Source reference.} Original lemma 066A: focused-seeds/AGP/sources/stacks-snapshot-20261011/tags.txt, source lines 7929--7929. Navigation pointer to a statement the counted proof refers to but does not contain; nothing load-bearing is delegated to it.\par\medskip
\ifdefined\hypertarget\hypertarget{sourcepointer-5}{}\fi
\noindent\textit{Source reference.} Original lemma 068M: focused-seeds/AGP/sources/stacks-snapshot-20261011/tags.txt, source lines 8011--8011. Navigation pointer to a statement the counted proof refers to but does not contain; nothing load-bearing is delegated to it.\par\medskip
\begin{lemma}
\label{lemma-Koszul-regular-flat-locally-regular}
Let $R$ be a ring. Let $f_1, \ldots, f_n$ be a Koszul-regular sequence
in $R$ such that $(f_1, \ldots, f_n) \not = R$.
Consider the faithfully flat, smooth ring map
$$
R \longrightarrow
S = R[\{t_{ij}\}_{i \leq j}, t_{11}^{-1}, t_{22}^{-1}, \ldots, t_{nn}^{-1}]
$$
For $1 \leq i \leq n$ set
$$
g_i = \sum\nolimits_{i \leq j} t_{ij} f_j \in S.
$$
Then $g_1, \ldots, g_n$ is a regular sequence in $S$ and
$(f_1, \ldots, f_n)S = (g_1, \ldots, g_n)$.
\end{lemma}

\begin{proof}
The equality of ideals is obvious as the matrix
$$
\left(
\begin{matrix}
t_{11} & t_{12} & t_{13} & \ldots \\
0 & t_{22} & t_{23} & \ldots \\
0 & 0 & t_{33} & \ldots \\
\ldots & \ldots & \ldots & \ldots
\end{matrix}
\right)
$$
is invertible in $S$.
Because $f_1, \ldots, f_n$ is a Koszul-regular sequence we see that
the kernel of
$R \to R^{\oplus n}$, $x \mapsto (xf_1, \ldots, xf_n)$ is zero (as it
computes the $n$the Koszul homology of $R$ w.r.t.\ $f_1, \ldots, f_n$).
Hence by
Lemma \ref{lemma-make-nonzero-divisor}
we see that $g_1 = f_1 t_{11} + \ldots + f_n t_{1n}$ is a nonzerodivisor
in $S' = R[t_{11}, t_{12}, \ldots, t_{1n}, t_{11}^{-1}]$. We see that
$g_1, f_2, \ldots, f_n$ is a Koszul-sequence in $S'$ by
Lemma \ref{lemma-koszul-regular-flat-base-change} and
\ref{lemma-independence-of-generators}.
We conclude that
$\overline{f}_2, \ldots, \overline{f}_n$ is a Koszul-regular sequence
in $S'/(g_1)$ by
Lemma \ref{lemma-truncate-koszul-regular}.
Hence by induction on $n$ we see that the images
$\overline{g}_2, \ldots, \overline{g}_n$ of $g_2, \ldots, g_n$ in
$S'/(g_1)[\{t_{ij}\}_{2 \leq i \leq j}, t_{22}^{-1}, \ldots, t_{nn}^{-1}]$
form a regular sequence. This in turn means that
$g_1, \ldots, g_n$ forms a regular sequence in $S$.
\end{proof}

\end{document}
